\section{The scheme-theoretic theorem of the cube}
\label{mumford-III10}
\dcref{M:ch3:III.10}

\begin{theorem}[Scheme-theoretic cube theorem]
\label{mumford-thm-cube-schemes}
\dcref{M:ch3:III.10}
Let $X$ and $Y$ be complete varieties, let $Z$ be a connected $k$-scheme with
a chosen section $z_0:\Spec k\to Z$, and let $L$ be a line bundle on
$X\times Y\times Z$.  If the restrictions to
\[
 \{x_0\}\times Y\times Z,\qquad X\times\{y_0\}\times Z,
 \qquad X\times Y\times\{z_0\}
\]
are trivial for chosen base points, then $L$ is trivial.  The proof uses the
maximal closed triviality locus and a nilpotent local-ring argument; the
conclusion is stable under base change that preserves the chosen section.
\uses{mumford-thm-cube-varieties,mumford-thm-seesaw}
\notready
\end{theorem}

\begin{proposition}[Theorem of the square]
\label{mumford-prop-square}
\dcref{M:ch3:III.10}
For an abelian variety $A$ and a line bundle $L$, the rule
\[
 \phi_L(x)=t_x^*L\otimes L^{-1}
\]
defines a homomorphism $\phi_L:A\to\Adual$.  Equivalently,
\[
 t_{x+y}^*L\otimes t_x^*L^{-1}\otimes t_y^*L^{-1}\otimes L
\simeq\cO_A
\]
functorially in $x,y$.
\uses{mumford-thm-cube-schemes,mumford-thm-dual-any-characteristic}
\notready
\end{proposition}

\begin{definition}[Biextension attached to a line bundle]
\label{mumford-def-biextension}
\dcref{M:ch3:III.10}
The line bundle
\[
 \mathcal B_L=m^*L\otimes p_1^*L^{-1}\otimes p_2^*L^{-1}
\]
on $A\times A$ is the biextension associated with $L$.  The cube theorem
gives its additivity in each variable and hence the homomorphism $\phi_L$.
\uses{mumford-prop-square}
\notready
\end{definition}

\section[Group schemes]{Basic theory of group schemes}
\label{mumford-III11}
\dcref{M:ch3:III.11}

\begin{definition}[Group scheme and $S$-points]
\label{mumford-def-group-scheme}
\dcref{M:ch3:III.11}
A group scheme over $k$ is a $k$-scheme $G$ with morphisms
$m:G\times G\to G$, $e:\Spec k\to G$, and $i:G\to G$ satisfying the group
identities as equalities of morphisms.  For every $k$-scheme $S$, the set
$G(S)=\Hom_k(S,G)$ is then a group functorially in $S$.
\notready
\end{definition}

\begin{lemma}[Kernels and subgroup schemes]
\label{mumford-lem-group-scheme-kernel}
\dcref{M:ch3:III.11}
The fibre over the identity of a homomorphism of group schemes is a closed
subgroup scheme.  Formation of this kernel commutes with base change, and a
homomorphism is a monomorphism precisely when its kernel is trivial as a
functor.
\uses{mumford-def-group-scheme}
\notready
\end{lemma}

\begin{proposition}[Invariant differentials and Lie algebra]
\label{mumford-prop-group-scheme-lie}
\dcref{M:ch3:III.11}
For a group scheme $G$ of finite type, translation identifies all tangent
spaces with $\Lie(G)=T_eG$.  Every tangent vector at the identity extends
uniquely to a left-invariant derivation; these derivations form the Lie algebra
of $G$, and a homomorphism induces a Lie-algebra map by its differential.  In
positive characteristic the differential alone does not determine the global
homomorphism (the Frobenius is the standard warning).
\uses{mumford-def-group-scheme,mumford-lem-group-scheme-kernel}
\notready
\end{proposition}

\begin{proposition}[Standard finite group schemes]
\label{mumford-prop-standard-finite-groups}
\dcref{M:ch3:III.11}
The group schemes $\Gm$, $\Ga$, $\mu_n=\Spec k[T]/(T^n-1)$, and
$\alpha_p=\Spec k[T]/(T^p)$ in characteristic $p$ are determined by their
Hopf algebras.  In characteristic $p$, $\mu_p$ and $\alpha_p$ are finite and
nonreduced, while $\Z/p\Z$ is etale.
\uses{mumford-def-group-scheme}
\notready
\end{proposition}

\begin{theorem}[Smoothness of an abelian variety]
\label{mumford-thm-abelian-smooth}
\dcref{M:ch3:III.11}
Every abelian variety is a smooth connected commutative group scheme.  Its
translation action makes the sheaf of differentials free, and
$\dim\Lie(A)=\dim A$.
\uses{mumford-def-abelian-variety,mumford-prop-group-scheme-lie,mumford-cor-abelian-commutative}
\notready
\end{theorem}

\begin{proposition}[Multiplication kernels in characteristic $p$]
\label{mumford-prop-p-rank}
\dcref{M:ch3:III.11}
For an abelian variety $A$ of dimension $g$ in characteristic $p>0$,
$A[p](\bar k)$ is isomorphic to $(\Z/p\Z)^f$ for an integer
$0\leq f\leq g$, called the $p$-rank.  The finite group scheme $A[p]$ has
rank $p^{2g}$ and carries a canonical connected--etale exact sequence.  A
product decomposition requires an additional splitting hypothesis.
\uses{mumford-prop-multiplication-n,mumford-thm-abelian-smooth}
\notready
\end{proposition}

\section{Quotients by finite group schemes}
\label{mumford-III12}
\dcref{M:ch3:III.12}

\begin{definition}[Finite group-scheme action]
\label{mumford-def-group-scheme-action}
\dcref{M:ch3:III.12}
An action of a finite group scheme $G$ on a scheme $X$ is a morphism
$G\times X\to X$ satisfying the unit and associativity diagrams.  It is free
if the graph map $G\times X\to X\times X$ is a monomorphism.
\uses{mumford-def-group-scheme}
\notready
\end{definition}

\begin{theorem}[Finite group-scheme quotient]
\label{mumford-thm-finite-groupscheme-quotient}
If a finite group scheme $G$ acts on a quasi-projective scheme $X$ with affine
orbits, the fppf quotient sheaf $X/G$ is represented by a scheme and
$X\to X/G$ is finite.  If the action is free, this map is a finite locally
free $G$-torsor; only in that torsor case does faithfully flat descent identify
quasi-coherent sheaves on the quotient with $G$-equivariant sheaves on $X$.
\dcref{M:ch3:III.12}
\uses{mumford-def-group-scheme-action,mumford-lem-group-scheme-kernel}
\notready
\end{theorem}

\begin{proposition}[Quotient of an abelian variety]
\label{mumford-prop-abelian-quotient}
\dcref{M:ch3:III.12}
If $H\subset A$ is a finite subgroup scheme, then the quotient $A/H$ is an
abelian variety and the quotient morphism $q:A\to A/H$ is a finite flat
homomorphism with kernel $H$.  Every finite flat homomorphism of abelian
varieties is obtained in this way.
\uses{mumford-thm-finite-groupscheme-quotient,mumford-thm-abelian-smooth}
\notready
\end{proposition}

\begin{theorem}[Isogeny factorization]
\label{mumford-thm-isogeny-factorization}
\dcref{M:ch3:III.12}
For an isogeny $f:A\to B$, there is an isogeny $g:B\to A$ and an integer
$n>0$ with $g\circ f=[n]_A$ and $f\circ g=[n]_B$.  Thus isogenies become
isomorphisms after tensoring Hom groups with $\Q$.
\uses{mumford-prop-abelian-quotient,mumford-prop-multiplication-n}
\notready
\end{theorem}

\section{The dual abelian variety in any characteristic}
\label{mumford-III13}
\dcref{M:ch3:III.13}

\begin{definition}[The stabilizer group scheme of a line bundle]
\label{mumford-def-KL}
\dcref{M:ch3:III.13}
For a line bundle $L$ on an abelian variety $A$, define the closed subgroup
scheme
\[
 K(L)(S)=\{x\in A(S)\mid t_x^*L_S\simeq L_S\}
\]
for every $k$-scheme $S$.
\uses{mumford-def-abelian-variety,mumford-lem-group-scheme-kernel}
\notready
\end{definition}

\begin{proposition}[The homomorphism attached to $L$]
\label{mumford-prop-phi-L}
\dcref{M:ch3:III.13}
The assignment $x\mapsto t_x^*L\otimes L^{-1}$ is a homomorphism
$\phi_L:A\to\Adual$ and its scheme-theoretic kernel is $K(L)$.
\uses{mumford-def-KL,mumford-prop-square}
\notready
\end{proposition}

\begin{theorem}[Representability of the dual]
\label{mumford-thm-dual-any-characteristic}
For every abelian variety $A$ over an algebraically closed field, the functor
of algebraically trivial line bundles rigidified at the origin is represented
by an abelian variety $\Adual$.  It carries a universal normalized Poincare line
bundle $\cP$ on $A\times\Adual$, and the construction is compatible with base
change.
\dcref{M:ch3:III.13}
\uses{mumford-def-KL,mumford-prop-abelian-quotient,mumford-thm-cube-schemes}
\notready
\end{theorem}

\begin{proposition}[Biduality in arbitrary characteristic]
\label{mumford-prop-biduality-any-characteristic}
\dcref{M:ch3:III.13}
The Poincare bundle gives a canonical isomorphism $A\simeq\widehat{\Adual}$;
for $f:A\to B$, the dual morphism $\widehat f:\Bdual\to\Adual$ is compatible
with composition, products, and base change.
\uses{mumford-thm-dual-any-characteristic}
\notready
\end{proposition}

\section[Cartier duality]{Duality of finite commutative group schemes}
\label{mumford-III14}
\dcref{M:ch3:III.14}

\begin{definition}[Cartier dual]
\label{mumford-def-cartier-dual}
\dcref{M:ch3:III.14}
For a finite commutative group scheme $G$, its Cartier dual is
$G^D=\underline{\Hom}(G,\Gm)$.  It represents multiplicative characters and
comes with a universal perfect pairing $G\times G^D\to\Gm$.
\notready
\end{definition}

\begin{theorem}[Cartier biduality]
\label{mumford-thm-cartier-biduality}
For finite commutative $G$, the evaluation map $G\to G^{DD}$ is an
isomorphism.  Cartier duality is exact on finite commutative group schemes and
interchanges etale and multiplicative-type parts.
\dcref{M:ch3:III.14}
\uses{mumford-def-cartier-dual}
\notready
\end{theorem}

\begin{proposition}[Weil pairing]
\label{mumford-prop-weil-pairing}
\dcref{M:ch3:III.14}
If $n$ is invertible in $k$, the commutative group scheme $A[n]$ has a perfect
alternating pairing
\[
 e_n:A[n]\times\Adual[n]\longrightarrow\mu_n
\]
induced by the Poincare bundle.  A polarization identifies the two factors and
produces the usual symplectic Weil pairing on $A[n]$.
\uses{mumford-thm-dual-any-characteristic,mumford-def-cartier-dual}
\notready
\end{proposition}

\section[Applications]{Applications to abelian varieties}
\label{mumford-III15}
\dcref{M:ch3:III.15}

\begin{definition}[Polarization]
\label{mumford-def-polarization}
\dcref{M:ch3:III.15}
A polarization of $A$ is a homomorphism $\lambda:A\to\Adual$ equal to
$\phi_L$ for an ample line bundle $L$.  It is principal when it is an
isomorphism.
\uses{mumford-prop-phi-L,mumford-thm-dual-any-characteristic}
\notready
\end{definition}

\begin{theorem}[Ample line bundles and finite kernels]
\label{mumford-thm-ample-finite-kernel}
\dcref{M:ch3:III.15}
For a line bundle $L$ on $A$, call $L$ nondegenerate when $\phi_L$ is an
isogeny (equivalently, $K(L)$ is finite).  An ample $L$ is nondegenerate, and
$\phi_L$ is a polarization precisely when it is induced by an ample $L$; the
sign matters, since $L^{-1}$ has the same finite kernel but is anti-ample.
\uses{mumford-def-polarization,mumford-def-KL,mumford-prop-abelian-quotient}
\notready
\end{theorem}

\begin{proposition}[Neron--Severi sequence]
\label{mumford-prop-pic-ns-sequence}
\dcref{M:ch3:III.15}
There is an exact sequence
\[
 0\longrightarrow\Piczero(A)\longrightarrow\Pic(A)\longrightarrow
 \NS(A)\longrightarrow0,
\]
where $\NS(A)$ is a finitely generated free abelian group.  Under
$L\mapsto\phi_L$, $\NS(A)$ embeds in $\Hom(A,\Adual)$.
\uses{mumford-thm-dual-any-characteristic,mumford-prop-phi-L}
\notready
\end{proposition}

\begin{corollary}[Dual isogeny]
\label{mumford-cor-dual-isogeny}
\dcref{M:ch3:III.15}
For an isogeny $f:A\to B$ and a polarization $\lambda_B$ on $B$, the pullback
polarization on $A$ is
\[
 \lambda_A=\widehat f\circ\lambda_B\circ f:A\longrightarrow\Adual.
\]
Relative to chosen polarizations one may form the Rosati adjoint
$f^\dagger=\lambda_A^{-1}\widehat f\lambda_B$.  The isogeny-factorization
theorem supplies a quasi-inverse with integral composites; the integer depends
on the chosen normalization and is not identified here with $\deg f$.
\uses{mumford-thm-isogeny-factorization,mumford-def-polarization,mumford-prop-biduality-any-characteristic}
\notready
\end{corollary}

\begin{remark}[Scheme-cube package]
\label{mumford3-cube-theorem-ii}
The second cube theorem is the scheme version with its triviality-locus and
nilpotent-thickening proof interface.
\dcref{M:ch3:III.10}
\uses{mumford-thm-cube-schemes,mumford3-cube-triviality-locus}
\notready
\end{remark}

\section{Supporting dependency interfaces}
\label{mumford-III-coverage-register}
The detailed scheme-theoretic constructions in this register support the
duality and arithmetic results.  Their proofs are treated as external inputs
where the available evidence is insufficient for reconstruction.

\begin{convention}[Scheme conventions]
\label{mumford3-standing-scheme-conventions}
The ground field is algebraically closed, schemes are of finite type over it,
and a point means a closed point unless an $S$-valued point is displayed.
\dcref{M:ch3:III.10}
\notready
\end{convention}

\begin{definition}[Maximal triviality locus]
\label{mumford3-cube-triviality-locus}
For $L$ on $X\times Y$ with $X$ complete, there is a maximal closed subscheme
$Y_L\subseteq Y$ such that the restriction of $L$ is pulled back from $Y_L$.
A map $S\to Y$ factors through $Y_L$ exactly when the corresponding family is
trivial along $X$ up to a bundle from $S$.
\dcref{M:ch3:III.10}
\uses{mumford-thm-seesaw}
\notready
\end{definition}

\begin{lemma}[Nilpotent step in the cube proof]
\label{mumford3-complete-connected-local-argument}
To prove that the triviality locus fills a connected base, pass to an Artinian
local thickening at a boundary point.  The first nonzero graded piece of its
ideal would define a class in the Kunneth decomposition of $H^1(X\times Y)$;
the three face trivializations force that class to vanish.
\dcref{M:ch3:III.10}
\uses{mumford3-cube-triviality-locus,mumford-thm-cube-schemes}
\notready
\end{lemma}

\begin{definition}[Functor of points]
\label{mumford3-functor-of-points}
A scheme $X$ determines the contravariant functor $S\mapsto\Hom(S,X)$.
Yoneda identifies natural transformations between representable functors with
unique morphisms of schemes.
\dcref{M:ch3:III.11}
\notready
\end{definition}

\begin{definition}[Translations and actions]
\label{mumford3-translations-and-actions}
An $S$-point of a group scheme acting on $X$ induces an automorphism of $X_S$;
translations compose according to the group law, and invariant morphisms are
constant on these functorial orbits.
\dcref{M:ch3:III.11}
\uses{mumford3-functor-of-points,mumford-def-group-scheme}
\notready
\end{definition}

\begin{definition}[Tangent vectors and invariant fields]
\label{mumford3-tangent-space-and-vector-field}
A tangent vector at $x$ is a derivation $\cO_{X,x}\to k$.  For a group scheme,
left-invariant vector fields are derivations compatible with multiplication.
\dcref{M:ch3:III.11}
\uses{mumford3-functor-of-points}
\notready
\end{definition}

\begin{proposition}[Invariant extension]
\label{mumford3-invariant-field-extension}
Every tangent vector at the identity extends uniquely to a left-invariant
vector field.  Evaluation at the identity therefore identifies such fields
with $\Lie(G)$.
\dcref{M:ch3:III.11}
\uses{mumford3-tangent-space-and-vector-field,mumford-prop-group-scheme-lie}
\notready
\end{proposition}

\begin{proposition}[Commutative Lie algebra]
\label{mumford3-commutative-lie-algebra}
For a commutative group scheme the bracket of invariant fields is zero.  In
characteristic $p$, the restricted $p$-operation retains infinitesimal data.
\dcref{M:ch3:III.11}
\uses{mumford3-invariant-field-extension}
\notready
\end{proposition}

\begin{theorem}[Characteristic-zero smoothness]
\label{mumford3-char-zero-smoothness}
Every finite-type group scheme over a characteristic-zero field is smooth and
reduced; invariant fields transport smoothness from the identity to all points.
\dcref{M:ch3:III.11}
\uses{mumford3-invariant-field-extension}
\notready
\end{theorem}

\begin{definition}[Infinitesimal subgroup neighbourhoods]
\label{mumford3-subgroup-kernel-neighborhoods}
The scheme-theoretic kernel of a homomorphism is the identity fibre.  In
characteristic $p$, suitable powers of the maximal ideal at the identity define
finite infinitesimal subgroup schemes with the same Lie algebra at high order.
\dcref{M:ch3:III.11}
\uses{mumford-lem-group-scheme-kernel,mumford-prop-group-scheme-lie}
\notready
\end{definition}

\begin{definition}[Hyperalgebra]
\label{mumford3-hyperalgebra}
The filtered duals of the local coordinate rings at the identity form a
convolution algebra of distributions.  Multiplication on the group supplies
the convolution product and ordinary multiplication supplies its coproduct.
\dcref{M:ch3:III.11}
\uses{mumford3-subgroup-kernel-neighborhoods}
\notready
\end{definition}

\begin{proposition}[Invariant differential operators]
\label{mumford3-invariant-differential-operators}
Each distribution in the hyperalgebra acts on functions as a left-invariant
differential operator, and convolution corresponds to composition.
\dcref{M:ch3:III.11}
\uses{mumford3-hyperalgebra,mumford3-invariant-field-extension}
\notready
\end{proposition}

\begin{proposition}[Free finite quotient]
\label{mumford3-free-quotient-flatness}
For a free action of a finite group scheme, $X\to X/G$ is finite locally free
of rank equal to the rank of $G$ and is a principal fppf torsor.
\dcref{M:ch3:III.12}
\uses{mumford-def-group-scheme-action,mumford-thm-finite-groupscheme-quotient}
\notready
\end{proposition}

\begin{definition}[Finite epimorphism]
\label{mumford3-finite-epimorphism}
A homomorphism of group schemes is an epimorphism here when it is faithfully
flat with the quotient universal property; in the finite setting its kernel is
a finite normal subgroup scheme.
\dcref{M:ch3:III.12}
\uses{mumford-lem-group-scheme-kernel}
\notready
\end{definition}

\begin{corollary}[Finite normal subgroups and quotients]
\label{mumford3-normal-finite-subgroup-correspondence}
Finite normal subgroup schemes correspond to finite quotient epimorphisms:
the subgroup is the scheme-theoretic kernel and the target represents the
quotient sheaf.
\dcref{M:ch3:III.12}
\uses{mumford3-finite-epimorphism,mumford-thm-finite-groupscheme-quotient}
\notready
\end{corollary}

\begin{definition}[Principal finite-group fibration]
\label{mumford3-principal-fibration}
A principal finite-group fibration is a finite locally free torsor
$U\to V$; its degree is the rank of the acting group scheme.
\dcref{M:ch3:III.12}
\uses{mumford3-free-quotient-flatness}
\notready
\end{definition}

\begin{theorem}[Euler characteristic of a finite torsor]
\label{mumford3-euler-characteristic-quotient}
For a principal finite-group fibration $\pi:U\to V$ and a coherent sheaf $F$
on $V$, one has
\[
 \chi(U,\pi^*F)=(\deg\pi)\chi(V,F).
\]
\dcref{M:ch3:III.12}
\uses{mumford3-principal-fibration,mumford-cor-flat-base-change}
\notready
\end{theorem}

\begin{remark}[Symmetric-group model of a finite map]
\label{mumford3-finite-morphism-symmetric-quotient}
Away from the collision locus, an ordered fibre product of a constant-degree
finite map is a principal symmetric-group cover; this construction reduces
several quotient calculations to the free case.
\dcref{M:ch3:III.12}
\uses{mumford3-principal-fibration}
\notready
\end{remark}

\begin{proposition}[Valued-point criterion for $K(L)$]
\label{mumford3-kernel-valued-point-criterion}
An $S$-point $x:S\to A$ factors through $K(L)$ exactly when
$t_x^*L_S\simeq L_S\otimes p_S^*N$ for some line bundle $N$ on $S$.
\dcref{M:ch3:III.13}
\uses{mumford-def-KL,mumford-thm-seesaw}
\notready
\end{proposition}

\begin{theorem}[Universal property of the dual]
\label{mumford3-poincare-universal-mapping}
A line bundle on $S\times A$, rigidified along $S\times\{0\}$ and fibrewise
algebraically trivial, is uniquely the pullback of the normalized Poincare
bundle along a morphism $S\to\Adual$.
\dcref{M:ch3:III.13}
\uses{mumford-thm-dual-any-characteristic,mumford-thm-seesaw}
\notready
\end{theorem}

\begin{theorem}[Cohomology of the Poincare bundle]
\label{mumford3-poincare-cohomology}
For $g=\dim A$, the normalized Poincare bundle on $\Adual\times A$ has a
single one-dimensional total cohomology contribution in degree $g$; equivalently
its derived pushforward is concentrated at the origin with the corresponding
shift.
\dcref{M:ch3:III.13}
\uses{mumford3-poincare-universal-mapping,mumford-thm-base-change-criterion}
\notready
\end{theorem}

\begin{corollary}[Structure-sheaf cohomology]
\label{mumford3-structure-sheaf-cohomology}
If $\dim A=g$, then $h^i(A,\cO_A)=\binom gi$ and the cohomology algebra is the
exterior algebra on $H^1(A,\cO_A)$.
\dcref{M:ch3:III.13}
\uses{mumford3-poincare-cohomology}
\notready
\end{corollary}

\begin{corollary}[Tangent space of the dual]
\label{mumford3-tangent-dual-h1}
There is a canonical identification
$\Lie(\Adual)\simeq H^1(A,\cO_A)$, compatible with pullback by homomorphisms.
\dcref{M:ch3:III.13}
\uses{mumford3-poincare-universal-mapping,mumford-prop-group-scheme-lie}
\notready
\end{corollary}

\begin{corollary}[Degree of the dual isogeny]
\label{mumford3-dual-isogeny-degree}
For an isogeny $f:A\to B$, the dual map $\widehat f:\Bdual\to\Adual$ is an
isogeny and $\deg\widehat f=\deg f$.
\dcref{M:ch3:III.13}
\uses{mumford-prop-biduality-any-characteristic,mumford-thm-isogeny-factorization}
\notready
\end{corollary}

\begin{proposition}[Divisorial duality criterion]
\label{mumford3-divisorial-correspondence}
A normalized line bundle on $A\times B$ determines mutually dual homomorphisms
$A\to\widehat B$ and $B\to\Adual$.  If either is an isomorphism, the bundle is
a Poincare bundle and $A,B$ are dual abelian varieties.
\dcref{M:ch3:III.13}
\uses{mumford-thm-cube-schemes,mumford3-poincare-universal-mapping}
\notready
\end{proposition}

\begin{theorem}[Representability of Cartier characters]
\label{mumford3-cartier-representability}
For finite commutative $G$, the functor
$S\mapsto\Hom_S(G_S,\Gm)$ is represented by the Cartier dual $G^D$, and the
evaluation pairing is universal.
\dcref{M:ch3:III.14}
\uses{mumford-def-cartier-dual}
\notready
\end{theorem}

\begin{example}[Basic Cartier duals]
\label{mumford3-cartier-examples}
Cartier duality exchanges the constant cyclic group $\Z/n\Z$ with $\mu_n$;
in characteristic $p$, it displays the different reduced and infinitesimal
behaviour of $\Z/p\Z$, $\mu_p$, and $\alpha_p$.
\dcref{M:ch3:III.14}
\uses{mumford-thm-cartier-biduality,mumford-prop-standard-finite-groups}
\notready
\end{example}

\begin{theorem}[Local and reduced components]
\label{mumford3-local-reduced-decomposition}
A finite commutative group scheme over a perfect field has canonical connected
and etale parts, and Cartier duality exchanges etale groups with groups of
multiplicative type.  The associated exact sequences need not split canonically.
\dcref{M:ch3:III.14}
\uses{mumford-thm-cartier-biduality}
\notready
\end{theorem}

\begin{proposition}[Primitive elements and Lie theory]
\label{mumford3-primitive-elements-lie}
Primitive elements in the dual Hopf algebra act as invariant derivations; this
identifies additive characters of the Cartier dual with $\Lie(G)$, compatibly
with the restricted $p$-operation.
\dcref{M:ch3:III.14}
\uses{mumford3-hyperalgebra,mumford-thm-cartier-biduality}
\notready
\end{proposition}

\begin{definition}[Height-one group scheme]
\label{mumford3-height-one-group}
In characteristic $p$, a finite group scheme has height one when it is supported
at the identity and every element of its maximal ideal has zero $p$th power.
\dcref{M:ch3:III.14}
\uses{mumford3-subgroup-kernel-neighborhoods}
\notready
\end{definition}

\begin{theorem}[Height-one groups and restricted Lie algebras]
\label{mumford3-height-one-equivalence}
Finite commutative height-one group schemes are classified by finite-dimensional
abelian restricted $p$-Lie algebras; subgroup schemes correspond to restricted
subalgebras.
\dcref{M:ch3:III.14}
\uses{mumford3-height-one-group,mumford3-primitive-elements-lie}
\notready
\end{theorem}

\begin{corollary}[Height-one actions]
\label{mumford3-height-one-action}
An action of a height-one commutative group scheme is equivalently a compatible
family of commuting derivations satisfying the corresponding restricted
$p$-power relations.
\dcref{M:ch3:III.14}
\uses{mumford3-height-one-equivalence,mumford3-invariant-differential-operators}
\notready
\end{corollary}

\begin{theorem}[Kernel of a dual isogeny]
\label{mumford3-dual-isogeny-kernel}
If $f:A\to B$ is an isogeny with kernel $K$, then
$\ker(\widehat f)$ is canonically Cartier dual to $K$.
\dcref{M:ch3:III.15}
\uses{mumford3-dual-isogeny-degree,mumford-thm-cartier-biduality}
\notready
\end{theorem}

\begin{definition}[Height-one isogeny]
\label{mumford3-height-one-isogeny}
An isogeny in characteristic $p$ is height one when its kernel is a height-one
finite group scheme.
\dcref{M:ch3:III.15}
\uses{mumford3-height-one-group,mumford-prop-abelian-quotient}
\notready
\end{definition}

\begin{theorem}[Height-one isogenies and $p$-Lie subalgebras]
\label{mumford3-height-one-isogeny-p-lie}
Height-one quotient isogenies of $A$ correspond to restricted $p$-Lie
subalgebras of $\Lie(A)$ through the Lie algebra of their kernels.
\dcref{M:ch3:III.15}
\uses{mumford3-height-one-equivalence,mumford3-height-one-isogeny}
\notready
\end{theorem}

\begin{theorem}[$p$-rank and Frobenius]
\label{mumford3-p-rank-lie-frobenius}
The $p$-rank is invariant under isogeny and duality and equals the dimension of
the semisimple part of Frobenius on the appropriate first cohomology/Lie data.
\dcref{M:ch3:III.15}
\uses{mumford-prop-p-rank,mumford3-tangent-dual-h1,
 mumford3-dual-isogeny-kernel}
\notready
\end{theorem}

\begin{theorem}[Restricted power and cohomological Frobenius]
\label{mumford3-lie-frobenius-compatibility}
Under $\Lie(\Adual)\simeq H^1(A,\cO_A)$, the restricted $p$-power operation is
compatible with the Frobenius action used to compute the $p$-rank.
\dcref{M:ch3:III.15}
\uses{mumford3-primitive-elements-lie,mumford3-tangent-dual-h1}
\notready
\end{theorem}

\begin{theorem}[Riemann--Roch and degree]
\label{mumford3-riemann-roch}
For $L=\cO_A(D)$ on a $g$-dimensional abelian variety,
\[
 \chi(L)=\frac{(D^g)}{g!},\qquad \deg(\phi_L)=\chi(L)^2.
\]
\dcref{M:ch3:III.16}
\uses{mumford-prop-euler-line-bundle,mumford-prop-phi-L}
\notready
\end{theorem}

\begin{definition}[Index of a nondegenerate line bundle]
\label{mumford3-index-definition}
A line bundle is nondegenerate when $\phi_L$ is an isogeny.  Its index $i(L)$
is the unique degree in which its coherent cohomology is nonzero.
\dcref{M:ch3:III.16}
\uses{mumford-thm-ample-finite-kernel,mumford-thm-line-bundle-index}
\notready
\end{definition}

\begin{theorem}[Real-root index polynomial]
\label{mumford3-index-polynomial-theorem}
For ample $H$ and nondegenerate $L$, the polynomial
$P(t)=\chi(H^{\otimes t}\otimes L)$ has real roots, counted with multiplicity,
and the sign convention recovers $i(L)$ from the number of positive roots.
\dcref{M:ch3:III.16}
\uses{mumford3-riemann-roch,mumford3-index-definition}
\notready
\end{theorem}

\begin{lemma}[Index growth bound]
\label{mumford3-index-growth-lemma}
Cohomology of tensor powers grows polynomially with the degree predicted by
Riemann--Roch; this bounds the possible jump of the index at a root of $P$.
\dcref{M:ch3:III.16}
\uses{mumford3-riemann-roch,mumford-thm-semicontinuity}
\notready
\end{lemma}

\begin{lemma}[Root-free invariance]
\label{mumford3-index-step-a}
Along a segment in the Neron--Severi space on which the Euler polynomial has no
root, the cohomological index is constant.
\dcref{M:ch3:III.16}
\uses{mumford3-index-polynomial-theorem,mumford-thm-semicontinuity}
\notready
\end{lemma}

\begin{lemma}[Two-variable index comparison]
\label{mumford3-index-step-b}
A root-free rectangle for the two-variable Euler polynomial identifies the
indices at its opposite endpoints, reducing the index theorem to one variable.
\dcref{M:ch3:III.16}
\uses{mumford3-index-step-a}
\notready
\end{lemma}

\begin{corollary}[Index under positive tensor powers]
\label{mumford3-index-power-corollary}
For every $n>0$, a nondegenerate $L$ satisfies $i(L^{\otimes n})=i(L)$.
\dcref{M:ch3:III.16}
\uses{mumford3-index-step-a}
\notready
\end{corollary}

\begin{proposition}[Subadditivity of index]
\label{mumford3-index-subadditivity}
When $L_1,L_2$, and $L_1\otimes L_2$ are nondegenerate,
$i(L_1\otimes L_2)\leq i(L_1)+i(L_2)$.
\dcref{M:ch3:III.16}
\uses{mumford3-index-step-b}
\notready
\end{proposition}

\begin{proposition}[Multiplicity bounds an index jump]
\label{mumford3-index-root-multiplicity-bound}
At a root of the index polynomial, the change in index is at most the
multiplicity of that root.
\dcref{M:ch3:III.16}
\uses{mumford3-index-growth-lemma,mumford3-index-step-b}
\notready
\end{proposition}

\begin{corollary}[Complex analytic index]
\label{mumford3-complex-index-corollary}
Over $\C$, the algebraic index is the number of negative eigenvalues of the
Hermitian form in the Appell--Humbert datum.
\dcref{M:ch3:III.16}
\uses{mumford3-index-polynomial-theorem,mumford-thm-appell-humbert}
\notready
\end{corollary}

\begin{lemma}[Separation of points]
\label{mumford3-support-separation}
Suitable translates of an ample divisor supply sections of a high tensor power
that separate any two distinct points of $A$.
\dcref{M:ch3:III.17}
\uses{mumford-thm-line-bundle-index,mumford-prop-square}
\notready
\end{lemma}

\begin{lemma}[Separation of tangent vectors]
\label{mumford3-tangent-separation}
High tensor powers of an ample line bundle separate nonzero tangent vectors.
In characteristic $p$, the proof uses the height-one subgroup generated by the
chosen vector.
\dcref{M:ch3:III.17}
\uses{mumford3-support-separation,mumford3-height-one-equivalence}
\notready
\end{lemma}

\begin{theorem}[Fourth and higher powers]
\label{mumford3-very-ample-power}
For an ample $L$, every $L^{\otimes n}$ with $n>3$ is very ample.
\dcref{M:ch3:III.17}
\uses{mumford3-support-separation,mumford3-tangent-separation,
 mumford-thm-high-power-very-ample}
\notready
\end{theorem}

\begin{remark}[Group-scheme structure package]
\label{mumford3-group-scheme}
The functorial group law recorded in the group-scheme definition is the
scheme-theoretic replacement for reasoning only with geometric points.
\dcref{M:ch3:III.11}
\uses{mumford-def-group-scheme,mumford3-functor-of-points}
\notready
\end{remark}

\begin{remark}[Restricted Lie package]
\label{mumford3-lie-algebra}
In characteristic $p$, $\Lie(G)$ carries its commutator bracket and the
restricted operation induced by $p$-fold composition of invariant derivations.
\dcref{M:ch3:III.11}
\uses{mumford-prop-group-scheme-lie,mumford3-invariant-field-extension}
\notready
\end{remark}

\begin{example}[Infinitesimal standard groups]
\label{mumford3-alpha-mu-examples}
The schemes $\alpha_p=\Spec k[T]/(T^p)$ and
$\mu_p=\Spec k[T]/(T^p-1)$ exhibit additive and multiplicative nonreduced
subgroups in characteristic $p$.
\dcref{M:ch3:III.11}
\uses{mumford-prop-standard-finite-groups}
\notready
\end{example}

\begin{remark}[Normal quotient package]
\label{mumford3-general-quotient-group-scheme}
Whenever the quotient by a normal subgroup scheme is representable, it inherits
a unique group law and is universal for invariant homomorphisms.
\dcref{M:ch3:III.11}
\uses{mumford3-normal-finite-subgroup-correspondence}
\notready
\end{remark}

\begin{remark}[Group-action package]
\label{mumford3-action}
Freeness is detected by the graph morphism $G\times X\to X\times X$, and
equivariant sheaves carry the cocycle needed for descent.
\dcref{M:ch3:III.12}
\uses{mumford-def-group-scheme-action}
\notready
\end{remark}

\begin{remark}[Finite quotient package]
\label{mumford3-finite-quotient-theorem}
The finite quotient theorem represents $X/G$ and separates representability
from the stronger free-torsor descent assertion.
\dcref{M:ch3:III.12}
\uses{mumford-thm-finite-groupscheme-quotient,mumford3-free-quotient-flatness}
\notready
\end{remark}

\begin{corollary}[Invariant pushforward]
\label{mumford3-invariant-pushforward}
For a finite quotient $\pi:X\to X/G$, the invariant subsheaf
$(\pi_*F)^G$ descends an equivariant $F$ in the free-torsor case.
\dcref{M:ch3:III.12}
\uses{mumford3-free-quotient-flatness}
\notready
\end{corollary}

\begin{remark}[Translation-kernel package]
\label{mumford3-kernel-line-bundle}
The functor $K(L)$ is represented by the scheme-theoretic kernel of $\phi_L$.
\dcref{M:ch3:III.13}
\uses{mumford-def-KL,mumford-prop-phi-L}
\notready
\end{remark}

\begin{remark}[Dual-construction package]
\label{mumford3-dual-construction}
The normalized Poincare bundle and the representing abelian variety together
constitute the construction of $\Adual$.
\dcref{M:ch3:III.13}
\uses{mumford-thm-dual-any-characteristic,mumford3-poincare-universal-mapping}
\notready
\end{remark}

\begin{remark}[$\phi_L$ package]
\label{mumford3-phi-l-morphism}
The theorem of the square makes $x\mapsto t_x^*L\otimes L^{-1}$ a homomorphism
whose kernel is $K(L)$.
\dcref{M:ch3:III.13}
\uses{mumford-prop-square,mumford-prop-phi-L}
\notready
\end{remark}

\begin{corollary}[Canonical biduality]
\label{mumford3-duality-hypothesis}
The Poincare correspondence identifies $A$ canonically with its bidual.
\dcref{M:ch3:III.13}
\uses{mumford-prop-biduality-any-characteristic,mumford3-divisorial-correspondence}
\notready
\end{corollary}

\begin{remark}[Cartier-duality package]
\label{mumford3-cartier-dual}
Cartier duality is the affine Hopf-algebra duality representing multiplicative
characters of finite commutative group schemes.
\dcref{M:ch3:III.14}
\uses{mumford-def-cartier-dual,mumford3-cartier-representability}
\notready
\end{remark}

\begin{definition}[$p$-rank]
\label{mumford3-p-rank}
The $p$-rank $f$ is determined by
$A[p](\bar k)\simeq(\Z/p\Z)^f$, with $0\leq f\leq\dim A$.
\dcref{M:ch3:III.15}
\uses{mumford-prop-p-rank}
\notready
\end{definition}

\begin{theorem}[Nondegenerate vanishing]
\label{mumford3-vanishing-theorem}
A nondegenerate line bundle has cohomology in one degree only; an ample bundle
has index zero.
\dcref{M:ch3:III.16}
\uses{mumford-thm-line-bundle-index,mumford3-index-definition}
\notready
\end{theorem}

\begin{corollary}[Structure-sheaf dimensions]
\label{mumford3-h0-structure-cohomology}
For $g=\dim A$, one has $h^i(A,\cO_A)=\binom gi$.
\dcref{M:ch3:III.16}
\uses{mumford3-structure-sheaf-cohomology}
\notready
\end{corollary}

\begin{corollary}[Projective realization]
\label{mumford3-projectivity-consequence}
A high very ample power of an ample line bundle realizes $A$ as a closed
projective subvariety.
\dcref{M:ch3:III.17}
\uses{mumford3-very-ample-power,mumford-cor-abelian-projective}
\notready
\end{corollary}

\section[Cohomology of line bundles]{Cohomology of line bundles}
\label{mumford-III16}
\dcref{M:ch3:III.16}

\begin{proposition}[Euler characteristic polynomial]
\label{mumford-prop-euler-line-bundle}
\dcref{M:ch3:III.16}
For a line bundle $L$ on a $g$-dimensional abelian variety, $\chi(L)$ depends
only on its Neron--Severi class and
\[
 \chi(L)=\frac{c_1(L)^g}{g!},\qquad \chi(L^{\otimes n})=n^g\chi(L).
\]
\uses{mumford-prop-pic-ns-sequence,mumford-thm-semicontinuity}
\notready
\end{proposition}

\begin{theorem}[Index and vanishing]
\label{mumford-thm-line-bundle-index}
\dcref{M:ch3:III.16}
If $L$ is nondegenerate, exactly one $H^i(A,L)$ is nonzero; its index is the
algebraic index attached to the numerical class (over $\C$ it is the number of
negative eigenvalues of the associated Hermitian form), and its
dimension is $|\chi(L)|$.  In particular, if $L$ is ample then
$H^i(A,L)=0$ for $i>0$ and $h^0(A,L)=\chi(L)>0$.
\uses{mumford-prop-euler-line-bundle,mumford-thm-ample-finite-kernel}
\notready
\end{theorem}

\begin{corollary}[Translation invariance of cohomology]
\label{mumford-cor-translation-cohomology}
\dcref{M:ch3:III.16}
For every $a\in A$, translation induces isomorphisms
$H^i(A,L)\simeq H^i(A,t_a^*L)$.  Cohomology dimensions are constant on the
translation orbit of a line bundle.
\uses{mumford-lem-abelian-translations,mumford-thm-line-bundle-index}
\notready
\end{corollary}

\section[Very ample line bundles]{Very ample line bundles}
\label{mumford-III17}
\dcref{M:ch3:III.17}

\begin{theorem}[High powers are very ample]
\label{mumford-thm-high-power-very-ample}
If $L$ is ample on an abelian variety, then $L^{\otimes n}$ is very ample for
every integer $n>3$.  The borderline $n=3$
requires the separate point- and tangent-separation argument and is not
asserted here.  The complete linear system defines a closed immersion
\[
 A\hookrightarrow\PP(H^0(A,L^{\otimes n})^\vee).
\]
\dcref{M:ch3:III.17}
\uses{mumford-thm-line-bundle-index,mumford-thm-cube-schemes}
\notready
\end{theorem}

\begin{corollary}[Projectivity and finite generation]
\label{mumford-cor-abelian-projective}
\dcref{M:ch3:III.17}
Every abelian variety is projective, and its homogeneous coordinate ring for a
high enough power of any ample line bundle is finitely generated.
\uses{mumford-thm-high-power-very-ample,mumford-def-abelian-variety}
\notready
\end{corollary}
